\exercice*

\begin{enumerate}

    (* for question in questions *)
    (* set q=questions[question] *)
  \item
  % (( question ))
    % (( q ))

    (* if question == "appliquer1" *)
    $(( q.a ))/(( q.b ))$ de \SI{(( q.c ))}{(( q.u ))} est égal à :
      \begin{align*}
        \frac{(( q.a ))}{(( q.b ))}\times \numprint{(( q.c ))}
        &=\frac{(( q.a ))\times  \numprint{(( q.c ))}}{(( q.b ))}\\
        &=\frac{(( q.a ))\times \numprint{(( q.c//q.b ))}\times\cancel{(( q.b ))}}{\cancel{(( q.b ))}}\\
        &=(( q.a ))\times \numprint{(( q.c//q.b ))}\\
        &=\numprint{(( q.a * q.c // q.b ))}
      \end{align*}
      La réponse est donc \SI{(( q.a * q.c // q.b ))}{(( q.u ))}.
    (* endif *)

    (* if question == "appliquer2" *)
    Le nombre de cadres dans l'entreprise est : $
  (( q.employes ))\times\frac{(( q.pourcent ))}{100}
  =
  \frac{(( q.employes )) \times (( q.pourcent ))}{100}
  =
  \frac{(( q.employes // 10 )) \times \cancel{10} \times (( q.pourcent // 10 )) \times \cancel{10}}{\cancel{10}\times\cancel{10}}
  =
  (( q.employes * q.pourcent // 100 ))
    $.

    (* endif *)

    (* if question == "inverse1" *)
    On peut résoudre cela avec un tableau de proportionnalité :
    \begin{center}
      \begin{tabular}{|l|c|c|}
        \hline
        Partie & (( q.partie )) & (( q.pourcent )) \\
        \hline
        Total  & (( q.lettre )) & 100\\
        \hline
      \end{tabular}
    \end{center}
    La valeur manquante (( q.lettre )) vaut donc :
    \begin{align*}
      (( q.lettre ))
      &= \frac{(( q.partie )) \times 100}{(( q.pourcent ))} \\
      &= \frac{(( q.partie * 10 // q.pourcent )) \times \cancel{(( q.pourcent // 10 ))} \times \cancel{10} \times 10}{\cancel{(( q.pourcent // 10))} \times \cancel{10}}\\
      &= (( q.total ))
    \end{align*}
    La quantité (( q.lettre )) vaut donc $(( q.total ))$.
    (* endif *)

    (* if question == "calculer1" *)
    La proportion de femmes dans l'entreprise est :
    $
    \frac{(( q.femmes ))}{(( q.employes ))}
    =\frac{(( q.femmes // q.pgcd ))\times\cancel{(( q.pgcd ))}}{(( q.employes // q.pgcd ))\times\cancel{(( q.pgcd ))}}
    =\frac{(( q.femmes // q.pgcd ))}{(( q.employes // q.pgcd ))}
    (* if 100*((q.femmes / q.employes)|round(2)) == 100 * (q.femmes / q.employes)  *)
    =(( (q.femmes / q.employes)|facteur() ))
    =(( (q.femmes / q.employes * 100)|facteur() ))\%
    (* else *)
    \approx (( (q.femmes / q.employes)|facteur("2") ))
    \approx (( (q.femmes / q.employes * 100)|facteur("0") ))\%
    (* endif *)
    $.
    (* endif *)

    (* if question == "calculer2" *)
    \begin{enumerate}
      \item Parmi (( q.notes.values()|sum )) élèves, (( q.notes[q.note] )) élèves ont eu (( q.note )) sur 5, soit :
        $\frac{(( q.notes[q.note] ))}{(( q.notes.values()|sum ))}=(( (q.notes[q.note]/(q.notes.values()|sum))|facteur() ))=(( (q.notes[q.note]/(q.notes.values()|sum)*100)|facteur() )) \%$.
      \item Parmi (( q.notes.values()|sum )) élèves, $
        (* for n in range(q.seuil, 6) *)
          (( q.notes[n] ))
          (* if not loop.last *)+(* endif *)
        (* endfor *)
        =
        (( q.audessusduseuil ))
        $ élèves ont eu plus de (( q.seuil )) sur 5, soit :
        $\frac{((q.audessusduseuil))}{(( q.notes.values()|sum ))}=(( (q.audessusduseuil/(q.notes.values()|sum))|facteur() ))=(( (q.audessusduseuil/(q.notes.values()|sum)*100)|facteur() )) \%$.
    \end{enumerate}
    (* endif *)

    (* if question == "calculer3" *)
    Les $\frac{(( q.a ))}{(( q.b ))}$ des $\frac{(( q.b ))}{(( q.c ))}$ de (( q.prix )) € donnent :
    \begin{align*}
      \frac{(( q.a ))}{(( q.b ))}\times\frac{(( q.b ))}{(( q.c ))}\times(( q.prix ))
      &= \frac{(( q.a )) \times \cancel{(( q.b ))} \times (( q.prix ))}{\cancel{(( q.b ))} \times (( q.c ))} \\
      &= (( (q.a * q.prix) // q.c ))
    \end{align*}
    Le résultat est donc (( (q.a * q.prix) // q.c )) €.
    (* endif *)

  (* endfor *)

\end{enumerate}
